\documentclass[11pt]{article}

\usepackage[margin=0.8in]{geometry}
\usepackage[T1]{fontenc}
\usepackage[utf8]{inputenc}
\usepackage{lmodern}
\usepackage[hidelinks]{hyperref}
\usepackage{xurl}
\usepackage{enumitem}
\usepackage{tabularx}
\usepackage{array}
\usepackage{titlesec}
\usepackage{xcolor}
\usepackage{needspace}

\setlength{\emergencystretch}{2em}
\setlength{\parindent}{0pt}
\setlength{\parskip}{2pt}
\pagestyle{empty}
\urlstyle{same}

\titleformat{\section}{\large\bfseries}{}{0em}{}[\titlerule]
\titlespacing*{\section}{0pt}{10pt}{6pt}

% ============================================================
% Helper commands
% ============================================================

\newcommand{\cvitem}[2]{%
\textbf{#1}: #2\\
}

\newcommand{\entry}[4]{%
\Needspace{6\baselineskip}
\begin{tabularx}{\textwidth}{@{}X>{\raggedleft\arraybackslash}p{0.24\textwidth}@{}}
\textbf{#1} & \textit{#2} \\
\multicolumn{2}{@{}l@{}}{\textit{#3}}
\end{tabularx}
\vspace{0.2em}
#4
\vspace{0.8em}
}

\newcommand{\subentry}[2]{%
\textit{#1} \hfill #2\\
}

\newcommand{\name}[1]{%
{\LARGE \textbf{#1}}\\[4pt]
}

\newcommand{\emaillink}[1]{%
\href{mailto:#1}{#1}%
}

\newcommand{\weblink}[1]{%
\href{https://#1}{#1}%
}

\newcommand{\githublink}[1]{%
\href{https://github.com/#1}{github.com/#1}%
}

\newcommand{\doilink}[1]{%
\href{https://doi.org/#1}{#1}%
}


\begin{document}

% ============================================================
% Header
% ============================================================

\begin{center}

{\Huge\bfseries Brian Kyanjo}\\[6pt]

{\large Computational Scientist \textbar\ Scientific Software Developer}\\[2pt]
{\small HPC \& Scientific Computing \textbar\ Numerical Modeling \textbar\ Open-Source Research Software}\\[8pt]

\small

\emaillink{bkyanjo3@gatech.edu}
\quad $\vert$ \quad
(208)~544--6609
\quad $\vert$ \quad
\href{https://github.com/KYANJO}{github.com/KYANJO}
\\[2pt]

\href{https://linkedin.com/in/kyanjo-brian}
{linkedin.com/in/kyanjo-brian}
\quad $\vert$ \quad
\href{https://orcid.org/0000-0002-0995-1051}
{ORCID: 0000-0002-0995-1051}

\end{center}

\vspace{4pt}
\hrule
\vspace{10pt}


% ============================================================
% Professional Summary
% ============================================================

\section*{Professional Summary}

Computational scientist and scientific software developer specializing in
high-performance scientific computing, numerical modeling, data assimilation,
and scalable research software. My work focuses on developing reliable,
reproducible, and interoperable computational systems for large-scale
scientific applications. I have experience developing open-source scientific
software, MPI- and GPU-based workflows, HPC and cloud computing platforms,
containerized environments, and coupled modeling systems across Python,
MATLAB, and compiled scientific codes.


% ============================================================
% Education
% ============================================================

\section*{Education}

\entry{Boise State University (BSU)}{2024}
{Computing Ph.D. in Computational Math Sciences and Engineering}
{
\textbf{Dissertation:}
\textit{GeoFlood: A Computational Model for Overland Flooding}\\
\textbf{Advisor:} Prof.\ Donna Calhoun
}

\entry{African Institute for Mathematical Sciences (AIMS), Rwanda}{2020}
{Master of Science in Mathematical Sciences (Climate Resilience)}
{
\textbf{Thesis:}
\textit{Simulating Tropical Cyclone over the Southwest Indian Ocean using
Model for Prediction Across Scales (MPAS): Sensitivity to Cloud Microphysics
Schemes}\\
\textbf{Advisor:} Prof.\ Babatunde J.\ Aboidun
}

\entry{Makerere University (MAK)}{2018}
{Bachelor of Science in Physical Sciences (Physics and Mathematics)}
{
\textbf{Final Project:}
\textit{Water Level Sensor System (Design and Construction)}\\
\textbf{Advisor:} Prof.\ Ireeta Tumps Winston
}


% ============================================================
% Technical Skills
% ============================================================

\section*{Technical Skills}

\cvitem{Scientific Computing}{Numerical PDEs, finite volume methods, adaptive mesh refinement (AMR), data assimilation, ensemble methods, uncertainty quantification, numerical modeling}

\cvitem{Programming}{Python, C++, C, FORTRAN, MATLAB, CUDA, R, Java, Bash}

\cvitem{HPC and Parallel Computing}{MPI, OpenMP, CUDA, Slurm, distributed-memory computing, GPU acceleration, parallel I/O, scalable scientific workflows}

\cvitem{Scientific Software Engineering}{Git, Linux, APIs, model coupling, Python--MATLAB interoperability, testing, workflow automation, reproducible computational workflows}

\cvitem{Research Infrastructure}{Docker, Apptainer, Spack, Jupyter, AWS/cloud computing, HPC deployment, containerized scientific environments}

\cvitem{Additional}{\LaTeX, SQL, HTML, electronics/auto technician experience}


% ============================================================
% Professional Experience
% ============================================================

\section*{Professional Experience}

\entry{Georgia Institute of Technology}{Aug.\ 2024 -- Present}
{Postdoctoral Research Fellow -- Computational Science and Scientific Software}
{
\begin{itemize}[leftmargin=1.2em, itemsep=2pt, topsep=2pt]

\item Design and develop ICESEE, an open-source, MPI-parallel ensemble data
assimilation framework for state and parameter estimation in ice-sheet and
geophysical models.

\item Develop scalable model-coupling interfaces and workflows for integrating
ICESEE with independent scientific models, including ISSM and Icepack.

\item Parallelize and optimize ensemble simulation and data assimilation
workflows for distributed-memory HPC systems using MPI and scalable
scientific computing methods.

\item Design and develop CryoStack, a scientific computing platform for
configuring, executing, monitoring, and managing reproducible scientific
workflows across institutional HPC and cloud computing environments.

\item Develop Remote and AWS cloud execution capabilities for scientific
workflows, including containerized execution on AWS Fargate and EC2 and
integration with institutional computing resources.

\item Develop reproducible scientific software environments using Docker,
Apptainer, Spack, and HPC deployment workflows.

\item Develop Python--MATLAB interoperability and ISSM--ICESEE coupling
workflows for ensemble forecasting, parameter estimation, and data
assimilation.

\item Develop workflow automation for experiment configuration, provenance,
simulation diagnostics, HPC job orchestration, and scientific result
management.

\item Contribute to research on scalable data assimilation and computational
methods for understanding ice-sheet dynamics and their contribution to
future sea-level change.

\end{itemize}
}

\entry{Georgia Institute of Technology}{Apr.\ 2025 -- Present}
{Research Faculty Advisory Council (RFAC), EAS}
{
\begin{itemize}[leftmargin=1.2em, itemsep=2pt, topsep=2pt]

\item Represent research faculty interests within the College of Sciences.

\item Serve as a liaison between research faculty and college leadership.

\item Advance initiatives supporting faculty engagement, inclusion, and
professional growth.

\item Evaluate internal grant proposals and contribute to funding decisions.

\end{itemize}
}


\entry{Rocky Mountain Advanced Computing Consortium (RMACC)}
{Oct.\ 2023 -- Sept.\ 2024}
{RMACC Apprenticeship}
{
\begin{itemize}[leftmargin=1.2em, itemsep=2pt, topsep=2pt]

\item Designed, quoted, and deployed a major expansion of the regional
high-performance computing cluster at the University of Colorado Boulder
(Alpine supercomputer).

\item Designed a research computing cluster in a Unix environment.

\item Supported HPC cluster provisioning, deployment automation, and systems
administration workflows.

\item Created playbooks and container images for MPI-based applications.

\end{itemize}
}


\entry{Boise State University, School of Computing}{June 2023 -- Dec.\ 2023}
{Graduate Teaching Assistant}
{
\begin{itemize}[leftmargin=1.2em, itemsep=2pt, topsep=2pt]

\item Developed a computational model for overland flooding:
\githublink{KYANJO/GeoFlood}.

\item Developed a CUDA-based version of GeoFlood.

\item Served as teaching assistant for Introduction to Programming in
Mathematics.

\end{itemize}
}


\entry{Boise State University, School of Computing}{Aug.\ 2020 -- July 2024}
{Graduate Research Assistant}
{
\begin{itemize}[leftmargin=1.2em, itemsep=2pt, topsep=2pt]

\item Developed GeoFlood, a computational model for overland flooding:
\githublink{KYANJO/GeoFlood}.

\item Benchmarked GeoFlood across CPU and GPU computing platforms.

\item Developed CUDA-based Riemann solvers for GeoFlood.

\item Developed ForestClaw extensions for atmospheric applications.

\item Explored data-driven and machine learning approaches to enhance
predictive modeling and simulation accuracy.

\item Implemented a discontinuous Galerkin solver within ForestClaw as a
summer research project.

\item Completed a comprehensive examination on Riemann solvers for wet/dry
interfaces in finite-volume schemes:
\href{https://github.com/KYANJO/Comp_exam/tree/main/code/comp_exam_artifact}
{github.com/KYANJO/Comp\_exam/tree/main/code/comp\_exam\_artifact}.

\item Implemented approximate and exact Riemann solvers for handling wet/dry
interfaces using finite-volume schemes.

\item Investigated accelerated wave-propagation algorithms using GPU and CPU
comparisons on HPC platforms.

\end{itemize}
}


\entry{Idaho National Laboratory (INL)}{May 2022 -- Sept.\ 2022}
{Intern}
{
\begin{itemize}[leftmargin=1.2em, itemsep=2pt, topsep=2pt]

\item Supported post-irradiation examination capabilities for nuclear fuel
using gamma tomography.

\item Simulated gamma-emission tomography of nuclear fuel and generated
tomographic reconstructions to identify isotopic distributions.

\item Performed thermal-fluid analyses of geothermal closed-loop
configurations and discussed results with the INL research team.

\item Participated in team meetings and contributed to maintenance and
improvement of the software suite.

\end{itemize}
}


\entry{AIMS Rwanda}{Sept.\ 2020 -- Dec.\ 2020}
{Tutor}
{
\begin{itemize}[leftmargin=1.2em, itemsep=2pt, topsep=2pt]

\item Tutored master's-level students in mathematics.

\item Assisted lecturers with academic quality assurance.

\item Supported administrative, library, IT, logistics, and technical
activities as needed.

\end{itemize}
}


\entry{ESMT Berlin}{2020}
{Certificate in Business Analytics}
{}


\entry{Marga Business Simulations, Germany}{2020}
{Certificate in Business Simulations Management}
{
\begin{itemize}[leftmargin=1.2em, itemsep=2pt, topsep=2pt]

\item Developed and marketed products within a competitive business simulation.

\item Collaborated with six students from different universities and
disciplines.

\item Served as team lead.

\end{itemize}
}


\entry{God Cares High School}{Jan.\ 2020 -- July 2020}
{Physics and Mathematics Teacher}
{
\begin{itemize}[leftmargin=1.2em, itemsep=2pt, topsep=2pt]

\item Taught physics and mathematics.

\item Supervised laboratory activities.

\end{itemize}
}


\entry{Ubuntu Hill School}{Aug.\ 2019 -- July 2020}
{Physics and Mathematics Teacher}
{
\begin{itemize}[leftmargin=1.2em, itemsep=2pt, topsep=2pt]

\item Taught physics and mathematics.

\end{itemize}
}


\entry{Kampala Capital City Authority (K.C.C.A.)}{May 2017 -- Aug.\ 2017}
{Intern}
{
\begin{itemize}[leftmargin=1.2em, itemsep=2pt, topsep=2pt]

\item Fabricated, repaired, and replaced damaged street lights and electronic
control boxes.

\item Troubleshot, maintained, and installed electric and solar street lights
across the city.

\item Supported maintenance of traffic lights and investigated approaches for
reducing peak-hour traffic.

\end{itemize}
}

\section*{Selected Funded Projects}

\begin{itemize}[leftmargin=1.2em, itemsep=3pt, topsep=2pt]

\item \textbf{Ar\^ete Glacier Initiative Sliding Law RFP},
\textit{A Framework for Detection of Changing Subglacial Properties through
Transient Data Assimilation}, 2026--2028.\\
Award: \$186,965. PI: Alexander Robel; Co-PI: Talea Mayo;
Co-Investigator: Brian Kyanjo. The project focuses on transient data
assimilation for detecting changes in subglacial properties and improving
constraints on future sea-level projections.

\end{itemize}
% ============================================================
% Open-Source Scientific Software
% ============================================================

\section*{Selected Scientific Software}

\begin{itemize}[leftmargin=1.2em, itemsep=4pt, topsep=2pt]

\item \textbf{ICESEE (Ice Sheet State and Parameter Estimator)} --
Open-source, MPI-parallel ensemble data assimilation framework for state
and parameter estimation in ice-sheet and geophysical models.\\
\githublink{ICESEE-project/ICESEE}

\item \textbf{CryoStack} --
Scientific computing platform for configuring, executing, monitoring,
and managing reproducible scientific workflows across institutional HPC
and AWS cloud environments.\\
\githublink{ICESEE-project/CryoStack}

\item \textbf{ICESEE-Containers} --
Reproducible container infrastructure for deploying MPI-based scientific
software across HPC systems using Docker and Apptainer.\\
\githublink{ICESEE-project/ICESEE-Containers}

\item \textbf{ICESEE-Spack} --
Spack-based software environment for reproducible builds and deployment
of ICESEE and coupled scientific modeling stacks.\\
\githublink{ICESEE-project/ICESEE-Spack}

\item \textbf{GeoFlood} --
Hybrid CPU/GPU computational model for overland flooding using finite
volume methods and adaptive mesh refinement.\\
\githublink{KYANJO/GeoFlood}

\item \textbf{ForestClaw Extensions} --
Numerical solver and adaptive mesh refinement development for
PDE-based geophysical simulations.\\
\githublink{KYANJO/forestclaw}

\end{itemize}


% ============================================================
% Publications
% ============================================================

\section*{Journal Publications and Refereed Conference Proceedings}

\begin{itemize}[leftmargin=1.2em, itemsep=4pt, topsep=2pt]

\item Brian Kyanjo and Alexander Robel.
\textit{The Ice Sheet State and Parameter Estimator (ICESEE) Library:
Ensemble Kalman Filtering for Ice Sheet Models}.
Submitted.
\doilink{10.48550/arXiv.2603.26947}.

\item Brian Kyanjo, Donna Calhoun, and David L.\ George.
\textit{A Hybrid CPU/GPU GeoFlood: High-Performance Computing in Overland
Flooding}.
In preparation.

\item Brian Kyanjo, Donna Calhoun, and David L.\ George.
\textit{GeoFlood: Computational Model for Overland Flooding}.
Submitted.
\doilink{10.48550/arXiv.2403.15435}.

\item B.\ Kyanjo, D.\ Calhoun, C.\ Burstedde, S.\ Aiton, J.\ Snively, and
M.\ Shih.
\textit{GPU Accelerated Adaptive Wave Propagation Algorithm}.
2023.
\doilink{10.13140/RG.2.2.11304.34569}.

\item Carlo Parisi, Paolo Balestra, Brian Kyanjo, Theron D.\ Marshall,
Travis L.\ McLing, and Mark D.\ White.
\textit{Closed Loop Geothermal Analysis Modeling and Simulation Using Idaho
National Laboratory's RELAP5-3D-FALCON Coupled Codes}.
Proceedings of the 48th Workshop on Geothermal Reservoir Engineering,
Stanford University, Stanford, California, February 6--8, 2023,
SGP-TR-224.

\end{itemize}


% ============================================================
% Conferences and Short Courses
% ============================================================

\section*{Selected Conferences and Short Courses}

\begin{itemize}[leftmargin=1.2em, itemsep=2pt, topsep=2pt]

\item American Geophysical Union (AGU) Annual Meeting, 2025

\item Society for Industrial and Applied Mathematics (SIAM),
5th Biennial Meeting of the SIAM Pacific Northwest Section, 2025

\item Scalable and Computationally Reproducible Approaches to Arctic Research,
2025

\item Adviser and Icepack Workshop and Hackathon, 2025

\item The International Conference for High Performance Computing,
Networking, Storage, and Analysis (SC), 2024

\item United States Research Software Engineers (US-RSE) Conference,
\textit{Software-Enabled Discovery and Beyond}, 2023

\item Accelerating Computing for Emerging Sciences (ACES) Workshop, 2023

\item SIAM Conference on Parallel Processing for Scientific Computing (PP22),
2022

\item SIAM Conference on Computational Science and Engineering,
attendee and poster judge

\item Winter School,
\textit{Advanced Numerical Methods for Hyperbolic PDE},
Laboratory of Applied Mathematics, University of Trento, Italy, 2021

\end{itemize}


% ============================================================
% Conference Posters
% ============================================================

\section*{Selected Conference Posters}

\begin{itemize}[leftmargin=1.2em, itemsep=2pt, topsep=2pt]

\item \textit{ICESEE: An Open-Source, Scalable Data Assimilation Library for
Ice Sheet Models}, AGU 2025

\item \textit{GeoFlood: A Computational Model for Overland Flooding},
2023 Graduate Showcase, Boise State University

\item \textit{GPU-Accelerated Wave Propagation Algorithm},
2023 Research Computing Days Conference

\item \textit{Geothermal Closed Loop},
2022 Clean Energy Development, Idaho National Laboratory

\item \textit{GPU-Accelerated Wave Propagation Algorithm},
2022 3rd Biennial Meeting of the SIAM Pacific Northwest Section

\end{itemize}


% ============================================================
% Invited Talks and Visits
% ============================================================

\section*{Invited Talks and Visits}

\begin{itemize}[leftmargin=1.2em, itemsep=2pt, topsep=2pt]

\item ESCO 2026:
\textit{Advanced Computational Methods for Earth System Modeling and Analysis}

\item \textit{ICESEE: An Open-Source, Scalable Data Assimilation Framework for
Ice Sheet Modeling}, 5th Biennial SIAM North Pacific Meeting

\item \textit{Math on ICE: ICESEE Coupling with ISSM for EnKF Data
Assimilation}

\item Seminar for the Geosystems Group:
\textit{Unified Computational Approaches to Flood and Ice Sheet Modeling}

\item \textit{Data Assimilation with Ice Sheet Modeling}, ART@SC24

\item \textit{GeoFlood: Computational Model for Overland Flooding},
4th Biennial Meeting of the SIAM Pacific Northwest Section

\item \textit{GeoFlood: A Computational Model for Overland Flooding},
2023 Lightning Talk, ACES Workshop

\item Research Visit at USGS:
\textit{Overland Flow Modeling with David L.\ George}, June 2023

\item Keynote Speaker, GAIN:
Promoting and encouraging study abroad for students from low-developing
countries

\end{itemize}


% ============================================================
% Activities, Associations, and Awards
% ============================================================

\section*{Activities, Associations, and Awards}

\begin{itemize}[leftmargin=1.2em, itemsep=2pt, topsep=2pt]

\item 5th Biennial SIAM North Pacific Meeting --
Mini-symposium organizer and poster judge,
\textit{Data Assimilation for Earth and Atmospheric Systems}
(October 2025)

\item Georgia Institute of Technology --
Research Faculty Advisory Council Representative for the School of Earth
\& Atmospheric Sciences (April 2025--Present)

\item Vanderbilt University --
Adviser \& Icepack Workshop and Hackathon Grand Prize Winner

\item Boise State University --
NSBE Public Policy SIG STEM Education, Team Co-Lead
(Oct.\ 2021 -- Oct.\ 2022)

\item Boise State University --
NSBE Student Chapter, President
(April 2023 -- April 2024)

\item Boise State University --
NSBE Student Chapter, Vice President
(April 2021 -- April 2023)

\item Boise State University --
SIAM Student Chapter, President
(Mar.\ 2022 -- April 2023)

\item 2023 Research Computing Days Hackathon --
Organizer and Judge

\item Idaho school outreach activities promoting STEM education

\item Boise State University --
SIAM Student Chapter, Vice President
(Mar.\ 2021 -- Mar.\ 2022)

\item 2021 SIAM Travel Award (AN21)

\item SIAM Travel Award (GS21)

\item Graduate Research Assistantship Scholarship (Ph.D.),
Aug.\ 2020 -- July 2024

\item 2023 Innovation Award Finalist,
Promoting STEM Across Idaho

\item Research Computing Days Conference --
Winning Poster,
\textit{GPU Wave Propagation Algorithm}

\item Association of State Floodplain Managers (ASFPM)

\item United States Research Software Engineers (US-RSE)

\item US-RSE 2023 Conference Travel Award

\item AIMS Rwanda --
AIMS Master's Scholarship (M.Sc.),
Aug.\ 2019 -- July 2020

\item ESMT Berlin --
Team Lead (May 2020 -- July 2020)

\item Makerere University Physics Society --
Vice President (April 2017 -- April 2018)

\item Government Scholarship on National Merit (B.Sc.),
Aug.\ 2015 -- May 2018

\end{itemize}


% ============================================================
% References
% ============================================================

\section*{References}

\begin{itemize}[leftmargin=1.2em, itemsep=3pt, topsep=2pt]

\item Prof.\ Alexander Robel,
Georgia Institute of Technology --
\emaillink{alexander.robel@eas.gatech.edu}

\item Prof.\ Donna Calhoun,
Boise State University --
\emaillink{donnacalhoun@boisestate.edu}

\item Mr.\ Craig Earley,
University of Colorado Boulder --
\emaillink{craig.Earley@colorado.edu}

\item Dr.\ David L.\ George,
U.S.\ Geological Survey --
\emaillink{dgeorge@usgs.gov}

\end{itemize}

\end{document}